\section{Study area and data}
\label{sec:data}

\subsection{Study area}
\label{subsec:area}

The study is carried out in the Izas experimental catchment, a small
(${\sim}0.55$\,km$^2$) alpine headwater basin located in the Central Spanish
Pyrenees (province of Huesca), at elevations of approximately
2000--2300\,m\,a.s.l. The catchment has been monitored for decades by the
Pyrenean Institute of Ecology (IPE--CSIC) and is characterised by a
strongly wind-influenced snowpack: persistent north-westerly winds generate
pronounced snow redistribution, with deep drifts on lee slopes and scoured
wind-exposed ridges. This produces the sharp, repeatable spatial patterns of
accumulation that make Izas an ideal testbed for high-resolution HS
prediction (Fig.~\ref{fig:study_area}). The LiDAR and meteorological data used
in this study were provided through a collaboration with IPE--CSIC.

\begin{figure}[htbp]
  \centering
  \includegraphics[width=\linewidth]{fig01_study_area.pdf}
  \caption{Location and physiography of the Izas experimental catchment
  (Central Spanish Pyrenees). (a) Regional setting within the Iberian
  Peninsula and the Pyrenees; (b) outline of the LiDAR-surveyed footprint
  (union of all acquisition dates) over a hillshade of the 1\,m DEM
  (EPSG:25830, UTM zone 30N); (c) example LiDAR-derived snow-depth map
  (1\,April 2025) illustrating the wind-driven accumulation patterns the
  models must reproduce.}
  \label{fig:study_area}
\end{figure}

\subsection{LiDAR snow-depth surveys}
\label{subsec:lidar}

Snow depth was derived from airborne LiDAR surveys by differencing snow-on
digital surface models against a snow-off reference, yielding HS maps at 1\,m
spatial resolution. A total of 27 acquisition dates spanning four hydrological
years (2021--2025) are available. The resulting HS rasters constitute the
prediction target (ground truth) of all models. To enable convolutional
processing, each dated HS map and its co-registered predictor stack are tiled
into fixed-size patches; pixels outside the catchment mask or without valid HS
are excluded from both training and evaluation.

\subsection{Predictor channels}
\label{subsec:channels}

Each training sample is a multi-channel image in which the predictors are
spatially co-registered with the target HS map at 1\,m resolution. The base
configuration comprises \textbf{22 channels}, which we group into four
families (Table~\ref{tab:channels}); Fig.~\ref{fig:channels} shows an example
predictor stack for one tile:

\begin{itemize}
  \item \textbf{Fine-scale topography (1\,m).} Elevation (DEM), slope,
  northness ($\cos$ of aspect) and eastness ($\sin$ of aspect), and the
  Topographic Position Index (TPI). These describe the local terrain that
  controls radiation, gravitational transport and small-scale deposition.

  \item \textbf{Coarse-scale topography (5\,m).} The same five terrain
  derivatives computed from a 5\,m DEM. The coarser scale captures the
  meso-topographic context (ridges, gullies, broad lee slopes) that governs
  wind redistribution over longer fetch distances.

  \item \textbf{Wind sheltering.} The maximum upwind slope index $S_x$
  \citep{winstral2002} computed over a 200\,m search distance for eight
  azimuths ($0^\circ$, $45^\circ$, \ldots, $315^\circ$). These eight channels
  encode directional exposure/sheltering and are the primary terrain-based
  proxy for wind-driven accumulation and scour.

  \item \textbf{Snow-cover history.} The fractional snow-cover persistence over
  the 15, 30 and 60 days preceding each acquisition, derived from a binary
  satellite snow-cover product (snow vs.\ no-snow) reprojected to 1\,m and
  temporally averaged; plus the instantaneous satellite snow-cover extent
  (SCE) of the acquisition date.
\end{itemize}

\begin{table}[t]
  \centering
  \caption{The 22 base predictor channels and the four additional
  meteorological channels (26-channel configuration). ``Static'' channels are
  constant in time; ``dynamic'' channels vary between acquisition dates.}
  \label{tab:channels}
  \small
  \begin{tabular}{clcc}
    \toprule
    Idx & Channel & Family & Type \\
    \midrule
    0  & DEM (1\,m)              & Topo 1\,m   & static  \\
    1  & Slope (1\,m)            & Topo 1\,m   & static  \\
    2  & Northness (1\,m)        & Topo 1\,m   & static  \\
    3  & Eastness (1\,m)         & Topo 1\,m   & static  \\
    4  & TPI (1\,m)              & Topo 1\,m   & static  \\
    5  & SCE (snow-cover extent) & Snow        & dynamic \\
    6--13 & $S_x$ 200\,m (8 azimuths) & Wind   & static  \\
    14 & Persistence 15\,d       & Snow        & dynamic \\
    15 & Persistence 30\,d       & Snow        & dynamic \\
    16 & Persistence 60\,d       & Snow        & dynamic \\
    17 & DEM (5\,m)              & Topo 5\,m   & static  \\
    18 & Slope (5\,m)            & Topo 5\,m   & static  \\
    19 & Northness (5\,m)        & Topo 5\,m   & static  \\
    20 & Eastness (5\,m)         & Topo 5\,m   & static  \\
    21 & TPI (5\,m)              & Topo 5\,m   & static  \\
    \midrule
    \multicolumn{4}{l}{\textit{Additional meteorological channels (26-ch
    configuration)}}\\
    22 & Mean air temp.\ 7\,d    & Meteo       & dynamic (scalar) \\
    23 & Mean air temp.\ 15\,d   & Meteo       & dynamic (scalar) \\
    24 & Mean air temp.\ 30\,d   & Meteo       & dynamic (scalar) \\
    25 & Accum.\ precip.\ from Oct & Meteo     & dynamic (scalar) \\
    \bottomrule
  \end{tabular}
\end{table}

\paragraph{Snow persistence.}
The persistence channels summarise the recent snow history of each pixel. For
a given acquisition date and temporal window $W\in\{15,30,60\}$ days, the
binary snow-cover product (snow $=1$, no-snow/cloud $=0$) is reprojected from
its native 10\,m resolution to 1\,m by bilinear resampling and averaged over
all available observations within the window, yielding a value in $[0,1]$ that
represents the fraction of recent days on which the pixel was snow-covered.
The ablation study (Section~\ref{subsec:res_ablation}) shows that these
channels are among the most informative predictors.

\paragraph{Meteorological channels (extended configuration).}
To test whether explicit meteorological forcing improves prediction, we build
an extended 26-channel configuration that adds four scalar channels derived
from the catchment automatic weather station (provided by IPE--CSIC): the mean
air temperature over the 7, 15 and 30 days preceding the acquisition, and the
precipitation accumulated since the 1st of October of the corresponding
hydrological year. These are scalar (spatially uniform) per date: every pixel
of a given acquisition shares the same value, so the channels encode the
temporal meteorological context rather than spatial structure.

\begin{figure}[htbp]
  \centering
  \includegraphics[width=\linewidth]{fig02_channels.pdf}
  \caption{Example predictor stack for a single $256\times256$\,m tile
  (2\,February 2021). Rows, top to bottom: fine 1\,m topographic derivatives
  (DEM, slope, northness, eastness, TPI); the corresponding coarse 5\,m
  derivatives; the eight $S_x$ wind-sheltering indices (0--315\textdegree);
  and the satellite snow-cover extent (SCE), the 15/30/60-day snow-persistence
  channels and the target LiDAR snow-depth map. Each panel is independently
  scaled; diverging fields are centred at zero.}
  \label{fig:channels}
\end{figure}

\subsection{Temporal split}
\label{subsec:split}

To obtain an honest estimate of predictive skill on \emph{unseen conditions},
we use a strictly temporal train/validation/test split rather than a random
pixel or tile split, which would leak spatial information between sets
(Fig.~\ref{fig:split}):

\begin{itemize}
  \item \textbf{Training:} hydrological years 2021--2023 (15 dates,
  4241 tiles).
  \item \textbf{Validation:} 2024 (7 dates, 431 tiles).
  \item \textbf{Test:} 2025 (5 dates, 220 tiles).
\end{itemize}

The test year 2025 was an exceptionally snowy season: the mean test HS
(1.09\,m) is $51\%$ higher than the mean training HS (0.72\,m). This
distribution shift is deliberate -- it stresses the extrapolation ability of
the models and explains the moderate absolute \Rsq{} values and the
non-negligible seed-to-seed variance observed in
Section~\ref{sec:results}. We regard robustness under this shift as a key,
realistic requirement for operational snow mapping.

\begin{figure}[htbp]
  \centering
  \includegraphics[width=\linewidth]{fig03_split_distribution.pdf}
  \caption{Temporal train/validation/test split and the corresponding
  distributions of LiDAR snow depth. (a) Per-date median snow depth
  ($\pm$ interquartile range) along the acquisition timeline, coloured by
  split. (b) Per-year snow-depth distributions (violins); black dots mark the
  annual means and the dashed line the training mean. The 2025 test year is
  markedly snowier than the 2021--2023 training years, illustrating the
  distribution shift that the models must extrapolate to.}
  \label{fig:split}
\end{figure}
